\section*{Chapter \ref{ch:iv3:cpp} --- C++}
\markboth{\protect\omTocMark{}{Answers: C++}}{\protect\omTocMark{}{Answers: C++}}

\begin{solution}{iq:iv3:cpp:1}
\texttt{p = \&b} and \texttt{*q = 3} compile, and the other four do not. Read a declaration from the name outward, to the left: \texttt{p} is a pointer to a \texttt{const int}, and \texttt{q} is a \texttt{const} pointer to an \texttt{int}.

So \texttt{p} can move to another object and cannot write through itself. \texttt{q} can write and cannot move, and \texttt{r} can do neither. \texttt{a} itself is not constant, so a read through \texttt{p} can still see a change made through \texttt{a} or \texttt{q}: \texttt{const int*} promises only that this path will not write.

\texttt{const int*} and \texttt{int const*} are the same type, since the position of \texttt{const} relative to \texttt{int} does not matter. Only its position relative to the \texttt{*} does.

\iqwrong{``\texttt{const int* p} means the value cannot change.'' It only stops writes through \texttt{p}.}
\iqfollow{And in a function signature?}{\texttt{void f(int* const p)} has the same type as \texttt{void f(int* p)}: the top-level \texttt{const} applies to the function's own copy of the argument, so it does not change the caller's view.}
\end{solution}

\begin{solution}{iq:iv3:cpp:2}
It prints \texttt{+global +static +auto -auto +auto -auto +heap end -static -global}. The heap object is never destroyed, since nothing calls \texttt{delete}, so \texttt{-heap} never appears.

\texttt{global} has static storage duration, and is built before \texttt{main}. \texttt{s} also has static storage duration, but is built the first time control passes through its declaration, which is why \texttt{+static} appears once. \texttt{a} has automatic storage, is built and destroyed on each call, and lives on the stack. \texttt{*h} has dynamic storage on the heap, and the pointer \texttt{h} itself is automatic.

At exit, static objects are destroyed in the reverse order of their construction, so \texttt{s} goes before \texttt{global}. The two static objects live in the program's data segment, not on the stack or the heap.

\iqwrong{``The operating system frees the heap object at exit, so its destructor runs.'' The memory is reclaimed, but no destructor runs. A destructor that flushes a file would never flush.}
\iqfollow{Is the static local thread-safe?}{Since C++11, its initialisation is: a second thread waits until it finishes.}
\end{solution}

\begin{solution}{iq:iv3:cpp:3}
Only default access differs: a \texttt{struct}'s members and bases are public by default, and a \texttt{class}'s are private. \texttt{S\{\}.x}, \texttt{C\{\}.get()} and \texttt{FromC\{\}.get()} compile, while \texttt{C\{\}.x} and \texttt{FromS\{\}.x} do not.

\texttt{FromC} is a \texttt{struct}, so it inherits publicly, and \texttt{get} stays public. \texttt{FromS} is a \texttt{class}, so it inherits privately, which hides \texttt{x} even though \texttt{x} is public in \texttt{S}. The default comes from the derived class, not from the base.

Both keywords can have methods, constructors, virtual functions and templates, and can be mixed freely in one hierarchy. Layout does not depend on the keyword.

The common convention is \texttt{struct} for a bundle of public data with no invariant, and \texttt{class} where members are kept private to protect one.

\iqwrong{``A \texttt{struct} is plain data, and it lives on the stack.'' That is C\#'s distinction, not C++'s. Where an object lives depends on how it is created.}
\end{solution}

\begin{solution}{iq:iv3:cpp:4}
It prints \texttt{2 3 2} and \texttt{2 4 4}. \texttt{r = b} assigns the value of \texttt{b} to \texttt{a}, because \texttt{r} is another name for \texttt{a}, and a reference cannot be reseated. \texttt{p = \&b} makes the pointer point elsewhere, and leaves \texttt{a} alone.

So after \texttt{r = b}, \texttt{a} is 2, and the later \texttt{b = 3} does not touch it. Then \texttt{*p = 4} writes to \texttt{b}.

Two other differences. A reference must be bound when it is declared, and it is never meant to be null, so a function taking one need not check. And a pointer is an object of its own, which you can take the address of, keep in a container and do arithmetic on, while \texttt{\&r} just gives the address of \texttt{a}.

\iqwrong{``A reference is safer, as it cannot dangle.'' It can. A reference to a local returned from a function dangles just as a pointer does.}
\iqfollow{When is a pointer the right parameter?}{When ``nothing'' is a valid argument, or when the function keeps it and may be pointed elsewhere later.}
\end{solution}

\begin{solution}{iq:iv3:cpp:5}
It prints 40, 8 and 10. Inside \texttt{main}, \texttt{a} is an array of 10 \texttt{int}, and \texttt{sizeof} gives 40 bytes. As a parameter, \texttt{int a[10]} means \texttt{int* a}, so the function sees a pointer, of 8 bytes, and the 10 in the brackets is ignored.

This decay is inherited from C: an array passed to a function turns into a pointer to its first element. The length is not passed at all, so the caller could hand over an array of 3, and the compiler would say nothing. g++ warns on this \texttt{sizeof} by default, which is why the listing has the warning switched off on a line not shown.

To keep the length, take a \texttt{std::span<int>} in C++20, which carries pointer and size and accepts arrays and vectors alike. Or take a reference to the array, \texttt{int (\&a)[10]}, or use a template over the length. \texttt{std::size} works on a true array because its type carries the length.

\iqwrong{``\texttt{sizeof(a) / sizeof(a[0])} gives the length.'' Inside the function it gives $8 / 4 = 2$, whatever was passed.}
\iqfollow{Why does \texttt{std::array} avoid this?}{It is a struct holding the array, so it passes by value or reference with its size in the type.}
\end{solution}

\begin{solution}{iq:iv3:cpp:6}
It prints \texttt{300 4}, then \texttt{44}, then \texttt{A 65 65}. \texttt{a + b} has type \texttt{int}, and nothing here is undefined.

Types smaller than \texttt{int} are promoted to \texttt{int} before arithmetic, so \texttt{a + b} is the integer 300, and \texttt{auto} makes \texttt{c} a 4-byte \texttt{int}. Storing 300 into \texttt{uint8\_t} reduces it modulo 256, which gives 44. That conversion to an unsigned type is defined. \texttt{-Wall -Wextra} does not warn, and nor does \texttt{-Wconversion} alone on g++ 11.4: it takes \texttt{-Wconversion -Warith-conversion}.

\texttt{std::uint8\_t} is an alias of \texttt{unsigned char} here, so \texttt{cout} prints it as a character, and 65 appears as \texttt{A}. \texttt{int(e)} or the unary plus, which promotes, prints the number.

The trap in feed parsers is that a one-byte field logged with \texttt{cout} prints as a control character, or as nothing.

\iqwrong{``\texttt{a + b} overflows, since both are 8 bits.'' The addition happens in \texttt{int}, so there is no overflow.}
\iqfollow{Can promotion cause undefined behaviour?}{Yes: two \texttt{uint16\_t} values of 65535 multiply as signed \texttt{int}, and that overflows.}
\end{solution}

\begin{solution}{iq:iv3:cpp:7}
It prints \texttt{3 2 3 1} and \texttt{1 3 0 3}. Parentheses choose among constructors as usual, so \texttt{a(3, 1)} is three copies of 1, and \texttt{c(3)} is three zeros. Braces prefer an \texttt{initializer\_list} constructor whenever one fits, so \texttt{b\{3, 1\}} is the two elements 3 and 1, and \texttt{d\{3\}} is one element.

The rule is that brace initialisation looks at \texttt{initializer\_list} constructors first, and considers the others only if none is viable. So for any container of numbers, braces mean ``these elements''.

With \texttt{std::vector<std::string> v\{3, \textquotedbl{}x\textquotedbl{}\}}, the list cannot convert, so the constructor taking a count and a value is used, and you get three strings. The same braces thus change meaning with the element type, which matters in templates.

\iqwrong{``Braces are the modern way, so use them everywhere.'' They also forbid narrowing, which is good, but on a container of numbers they do something other than parentheses.}
\iqfollow{What is a good habit?}{Use parentheses for a size and braces for a list of elements, and do not mix them in generic code without thinking.}
\end{solution}

\begin{solution}{iq:iv3:cpp:8}
It prints \texttt{+a +b body -b -a}. Members are initialised in the order they are declared in the class, whatever the order in the initialiser list, and destroyed in the reverse order. \texttt{-Wreorder}, which \texttt{-Wall} includes, reports the mismatch.

The listing turns off that warning on a line not shown, as it would otherwise fail a build with \texttt{-Werror}.

It matters when one initialiser depends on another. With \texttt{Pair() : b(a.name), a(\textquotedbl{}a\textquotedbl{})}, \texttt{b} is built second, so that one happens to work. The reverse dependency is the bug: if \texttt{a}'s initialiser reads \texttt{b}, it reads an object that does not yet exist, which is undefined, and reading the list suggests otherwise.

The reason for a fixed order is destruction. There may be several constructors with different lists, but only one destructor, and it must undo the construction in reverse.

\iqwrong{``The list decides, as it runs left to right.'' Base classes come first, and then the members in declaration order.}
\iqfollow{What of the constructor body?}{It runs after every member is built, so an assignment in the body is a second write, not an initialisation.}
\end{solution}

\begin{solution}{iq:iv3:cpp:9}
It prints \texttt{int}, \texttt{pointer}, \texttt{pointer}, \texttt{pointer}, \texttt{int}. \texttt{nullptr} has type \texttt{std::nullptr\_t}, which converts to any pointer type and not to \texttt{int}. The literal 0 is an \texttt{int}, so it picks the \texttt{int} overload.

Before C++11, null was written 0 or \texttt{NULL}, and either could silently choose an overload for integers. A literal 0 also stops being a null pointer once it passes through a template: in \texttt{relay(0)}, \texttt{T} is \texttt{int}, and an \texttt{int} variable holding 0 cannot convert to a pointer, so if only the pointer overload existed, \texttt{relay(0)} would not compile. \texttt{relay(nullptr)} keeps the type, and that still converts.

So \texttt{nullptr} says ``no object'' in the type system, and it survives being passed around. Its value is not the issue: on x86-64 Linux the null pointer is the all-zero address, but no code should depend on that.

\iqwrong{``0 and \texttt{nullptr} are the same thing.'' In \texttt{f(0)} they would call different functions, which is the whole point.}
\end{solution}

\begin{solution}{iq:iv3:cpp:10}
It prints \texttt{5: 1 2 3 0 3} and then \texttt{3}. \texttt{std::remove} works on iterators, so it can only move elements around. It cannot change a container's size, as it never sees the container. In C++20, \texttt{std::erase(v, 0)} does both steps.

The kept elements are shifted to the front, and the function returns the new logical end. What lies beyond is valid but unspecified, which is why the tail here reads \texttt{0 3}: those are the old positions, left as they were. For a type with a move, they may be moved-from.

The classic form is \texttt{v.erase(std::remove(v.begin(), v.end(), 0), v.end())}, which is linear, with each kept element moved at most once. That is far better than erasing one element at a time, which is quadratic.

\iqwrong{Using \texttt{v.end()} after \texttt{std::remove} as though it had moved. It has not, so a loop to \texttt{v.end()} sees the leftovers.}
\iqfollow{What about a \texttt{std::list}?}{It has its own member \texttt{remove}, which unlinks nodes and does change the size.}
\end{solution}

\begin{solution}{iq:iv3:cpp:11}
It prints 0 and then 1.8. In \texttt{accumulate}, the type of the running total is the type of the initial value, so a literal 0 makes it an \texttt{int}, and each step truncates.

So the steps are $0 + 0.6 = 0.6$, stored in an \texttt{int} as 0, three times over. The result is 0, and not even 1 or 2, since nothing ever carries. With fills of 0.6 the total stays 0 forever, which on a position would be serious.

\texttt{-Wall -Wextra} say nothing, since the conversion occurs inside a library template, in a system header, where warnings are off by default. Write the start as \texttt{0.0}, or use \texttt{std::reduce} with a typed start, or set the type yourself: \texttt{accumulate(b, e, double\{\})}.

The same trap applies with a start of \texttt{0} on a sum of \texttt{long} values, where each step's \texttt{long} result is cut down to \texttt{int}, so a large total comes out silently wrong.

\iqwrong{``It rounds, so the answer is near 2.'' Each step cuts toward zero and starts again from what remains.}
\end{solution}

\begin{solution}{iq:iv3:cpp:12}
It prints \texttt{1 2}, \texttt{10 20} and \texttt{10 7}. The first loop changes a copy, and the second changes the elements. \texttt{auto} leaves out the reference and the top-level \texttt{const}, so \texttt{a} is a plain \texttt{int}, and writing to it leaves \texttt{v} alone.

To keep them, say so: \texttt{auto\& b = r} deduces \texttt{const int\&} and \texttt{const auto\&} says it outright, and both bind to the element. \texttt{decltype(auto)} keeps the exact type of the initialiser.

For a range-based loop, use \texttt{const auto\&} on anything that may be costly to copy, such as a string or a large struct, and \texttt{auto\&} to change elements. Plain \texttt{auto} is right for small values, such as an integer or a \texttt{string\_view}, where a copy is cheap and the code sees no aliasing.

\iqwrong{\texttt{for (auto x : orders) x.qty = 0;} ``It compiled, so it worked.'' It zeroed copies, with no warning, and the orders are unchanged.}
\iqfollow{Why can \texttt{const auto\&} go wrong on a map?}{With a mismatched element type, such as \texttt{const std::pair<int, double>\&} for a map's \texttt{pair<const int, double>}, each element is converted to a temporary and bound: an unnoticed copy per step.}
\end{solution}

\begin{solution}{iq:iv3:cpp:13}
\texttt{b = a} copies the pointer, so both objects own one array. When \texttt{b} dies, the array is freed, and AddressSanitizer reports a \texttt{heap-use-after-free} at \texttt{a.data[0] = 2.0}. If that line were gone, \texttt{a}'s destructor would free the array a second time.

The compiler writes the copy operations itself unless told otherwise, and it copies members one by one, which for a raw pointer means the address, not the memory it points at.

The rule of five: if a class needs a custom destructor, copy constructor, copy assignment, move constructor or move assignment, it probably needs to think about all five. \Cref{lst:iv3:cp:bufferfive} does that. Assignment by value serves both copying and moving, and its swap means self-assignment cannot go wrong. \texttt{noexcept} on the moves lets containers use them.
\omcode{code/interviews-3/03-cpp/cpp/buffer_five.cpp}{1}{24}{All five, with assignment by value.}\label{lst:iv3:cp:bufferfive}

Better still is the rule of zero: hold a \texttt{std::vector<double>}, and write none of them. The compiler's versions are then correct, since each member manages itself. Writing the five is for the few classes whose whole job is to own one resource.

\iqwrong{``Declare the copy constructor private and be done.'' That stops the copying, and in modern C++ you would write \texttt{= delete}; it is fine for a type that should not be copied, but you still need to decide about moves.}
\end{solution}

\begin{solution}{iq:iv3:cpp:14}
It prints \texttt{0 1} and then \texttt{1 0}. A user-declared destructor means the compiler does not generate the move constructor or move assignment, so \texttt{std::move} falls back on the copy constructor. The program is still correct, and has lost only speed. The fix is the line \texttt{Logged(Logged\&\&) = default;}.

That one line is not quite all. Declaring a move constructor deletes the implicit copy operations, so \texttt{Logged(const Logged\&) = default;} and the two assignments should follow, if they are wanted, and \texttt{Logged() = default;} too.

Nothing warns you, as this is a silent performance change: \texttt{std::move} still compiles, and a copy is a valid way to meet a request to move. So a member like a long string or a big vector is now duplicated on every move.

The lesson is about what looks harmless. A destructor added to log something looks free, and it can change how every object of the class moves. If you need to see moves, count them, as the listing does.

\iqwrong{``An empty destructor is the same as none.'' Its body does nothing, but it is user-declared, and that is what the language looks at: the moves go, and a type that was trivially copyable stops being so, which also changes how it is passed.}
\iqfollow{What does \texttt{= default} on the destructor do?}{It is still user-declared, so it suppresses the moves in just the same way.}
\end{solution}

\begin{solution}{iq:iv3:cpp:15}
It prints \texttt{0 copies, 7 moves}, then \texttt{7 copies, 0 moves}. The capacity goes 1, 2, 4, then 8, so the five insertions move 1, 2 and 4 old elements. A vector moves them only if the move cannot throw. Otherwise it copies, to keep the strong exception guarantee of \texttt{push\_back}.

The guarantee says that if the insertion throws, the vector is as it was. With copies, that is easy to keep: if the third copy throws, destroy the new copies, free the new block, and the old elements are untouched. With a throwing move, the first elements of the old block are already moved-from when the third move throws, and moving them back could throw as well, so there is no way back.

So the vector uses \texttt{std::move\_if\_noexcept}, which gives a move only if the move is \texttt{noexcept} or the type cannot be copied at all. A move-only type with a throwing move still moves, and the guarantee is given up.

So a move constructor that never throws should say \texttt{noexcept}. Without it, the moves you wrote may never run in a vector. A compiler-generated move is \texttt{noexcept} on its own when all the members' moves are.

\iqwrong{``The compiler can see that this move cannot throw.'' It checks the declaration. An inline body that does nothing still counts as potentially throwing here.}
\iqfollow{How do you check a type?}{With a \texttt{static\_assert} next to the type, on the nothrow-move trait from \texttt{<type\_traits>}.}
\end{solution}

\begin{solution}{iq:iv3:cpp:16}
It prints \texttt{0 40}, \texttt{3}, \texttt{1 1}, \texttt{1 5} and \texttt{new!}. The standard promises the 40, the 3, the \texttt{1 1}, the \texttt{1 5} and the last line. The size 0 of a moved-from string is what libstdc++ does here, and the standard calls it only ``valid but unspecified''.

\texttt{std::move} does nothing at run time. It is a cast to an rvalue reference, which lets a move constructor be chosen. The last two lines of the listing show that: \texttt{r} is just another name for \texttt{a}, and \texttt{a} changes through it, with nothing moved at all.

The rest follows from what is moved. An \texttt{int} has no resources, so its move is a copy, and \texttt{i} stays 1. A \texttt{unique\_ptr} is promised to be null after being moved from, as its specification says so.

On a moved-from string, you may do whatever has no preconditions: assign to it, as the second line shows, destroy it, or ask for its size. You may not assume a value, so \texttt{a.front()} is not safe unless you check.

\iqwrong{``After \texttt{std::move(a)}, \texttt{a} is empty.'' The cast alone changes nothing, and even a real move does not promise an empty string.}
\iqfollow{What does a short string do?}{Its characters are in the object itself, so the move copies them. Nothing is cheaper there.}
\end{solution}

\begin{solution}{iq:iv3:cpp:17}
It prints \texttt{plain: 0 copies, 0 moves}, \texttt{forced: 0 copies, 1 moves} and \texttt{either: 1 copies, 0 moves}. \texttt{std::move} on the return prevents elision, as the returned expression is no longer the plain local name. In \texttt{either}, the returned expression \texttt{c ? p : q} is not the name of a local, so the result is copied.

Returning a local by name lets the compiler build it directly in the caller's slot, which g++ does here. That is allowed and not required. If it declines, the language still tries a move first, so \texttt{return t;} never needs \texttt{std::move}. So the cast gains nothing and costs a move. \texttt{-Wpessimizing-move}, in \texttt{-Wall}, warns about it, and the listing has it turned off on a line not shown.

The conditional is a different case. It is not a name, so no implicit move applies. \texttt{if (c) return p; return q;} would move, or could elide.

\texttt{make} works for a different reason. Since C++17, a function returning a prvalue directly has no temporary to copy at all: \texttt{Pinned\{\}} initialises \texttt{p} itself. So this is a rule of the language, and no copy constructor need exist.

\iqwrong{``Always write \texttt{std::move} on the return, to be sure it moves.'' That is the one place where it is harmful. It has a use only when returning a member, or a parameter that is an rvalue reference, in older standards.}
\end{solution}

\begin{solution}{iq:iv3:cpp:18}
It prints \texttt{copy}, \texttt{copy}, \texttt{copy}, \texttt{move}, \texttt{move}. Inside \texttt{bad}, \texttt{x} is a named variable, and any named variable is an lvalue, so \texttt{take(x)} always picks the \texttt{const std::string\&} overload.

\texttt{std::forward<T>(x)} restores what the caller passed: an lvalue stays an lvalue, and an rvalue becomes an rvalue again. It can do this because \texttt{T} carries that information: for a caller's lvalue, \texttt{T} is \texttt{std::string\&}, and for an rvalue it is \texttt{std::string}.

\texttt{std::move} makes everything an rvalue, without condition. With it in \texttt{good}, the call \texttt{good(s)} would move out of the caller's \texttt{s}, emptying a variable that the caller never gave up. That is a nasty bug, since the caller's code shows no \texttt{std::move}.

So the rule is simple: \texttt{std::forward} goes with a forwarding reference, and \texttt{std::move} with an rvalue reference. Forward an argument once at most, as a second use reads a moved-from object.

\iqwrong{``\texttt{x} has type \texttt{T\&\&}, so it is an rvalue.'' The type of the variable and the category of the expression are different things. A name is always an lvalue.}
\iqfollow{What goes wrong with a function pointer or a brace list?}{Perfect forwarding is not perfect. \texttt{good(\{1, 2\})} fails, as a brace list has no type from which to deduce \texttt{T}. A function pointer forwards fine, but the name of an overloaded function does not: it too has no single type from which to deduce \texttt{T}.}
\end{solution}

\begin{solution}{iq:iv3:cpp:19}
It prints \texttt{T is a reference}, \texttt{T is a value} and \texttt{T is a reference}. Only \texttt{T\&\&} in \texttt{f} and \texttt{auto\&\&} are forwarding references, as each has the exact form of a deduced name plus \texttt{\&\&}. \texttt{g(l)} does not compile, since \texttt{std::vector<T>\&\&} is a plain rvalue reference and \texttt{l} is an lvalue.

With \texttt{f(i)}, \texttt{i} is an lvalue, so a special deduction rule makes \texttt{T} equal to \texttt{int\&}. The parameter type becomes \texttt{int\& \&\&}. You cannot write that yourself, but it may arise from a template or an alias, and reference collapsing reduces it: if either part is \texttt{\&}, the result is \texttt{\&}, and only \texttt{\&\& \&\&} gives \texttt{\&\&}. So the parameter is \texttt{int\&}. With \texttt{f(1)}, \texttt{T} is \texttt{int}, and the parameter is \texttt{int\&\&}. With a \texttt{const int}, \texttt{T} is \texttt{const int\&}.

Decoration of any kind means it is no longer a forwarding reference, as in \texttt{const T\&\&} or in \texttt{std::vector<T>\&\&}. A \texttt{T\&\&} in a member of a class template, where \texttt{T} was fixed when the class was instantiated, is not one either, because no deduction happens at that call.

So \texttt{auto\&\& a = i} makes \texttt{a} an \texttt{int\&}, and \texttt{auto\&\& b = 1} an \texttt{int\&\&}, which the listing checks with \texttt{static\_assert}.

\iqwrong{``\texttt{\&\&} always means an rvalue reference.'' Only where nothing is deduced.}
\end{solution}

\begin{solution}{iq:iv3:cpp:20}
\texttt{make()} returns a temporary, which is destroyed at the end of the declaration, so \texttt{v} points at freed memory. AddressSanitizer reports \texttt{heap-use-after-free} at the line that reads \texttt{v}. A 5-character string would be held inside the \texttt{std::string} object, by the small-string optimisation, so no heap is involved, and the view would dangle into dead stack instead.

A \texttt{string\_view} is a pointer and a length with no ownership. Binding a \texttt{const std::string\&} to a temporary extends the temporary's life. A view has no such effect, as its conversion runs through an operator, and the temporary dies as usual. So keep the string in a named variable, and take the view from that.

As a parameter, a view is excellent: \texttt{void parse(std::string\_view)} accepts literals and strings with no copy, and the caller's argument outlives the call. As a member, or a return value, it is poor, since the object can outlive what it views, and nothing in the type says so. If a class needs to keep the text, it should hold a \texttt{std::string}.

\iqwrong{``It prints the right thing, so it is fine.'' Freed memory often still holds the old bytes. The program is undefined, and it passes until an allocation reuses the block.}
\iqfollow{Is the same true of a \texttt{std::span}?}{Yes. It is a pointer and a size, and has the same lifetime rule.}
\end{solution}

\begin{solution}{iq:iv3:cpp:21}
It prints \texttt{4 12}, then \texttt{4}. \texttt{strlen} reads until a zero byte. \texttt{sym.data()} points into the whole literal, and the only zero byte is at its end, so the C function sees \texttt{ESZ6,4501.25}, 12 characters, while the view has only 4.

So the bug is an assumption: a view is a pointer plus a length, and no zero byte need follow it. A C function has no way to receive the length, so it reads on. The substring of a longer text ends in the middle of other characters, and that is the usual case for a parser.

It is not undefined here, as the literal really does have a zero byte at the end. It is a wrong answer. With a view into a non-terminated buffer, such as bytes off a socket, \texttt{strlen} would run off the end, which is undefined.

It passes tests because a view of an entire literal does end at a zero byte, so the simple case works. Only a true substring shows it. The fix for a C call is to make a terminated copy, \texttt{std::string(sym).c\_str()}, as in the second line, and to keep that string alive through the call. Or use a function that takes a length.

\iqwrong{``\texttt{string\_view} has \texttt{c\_str()}, surely.'' It has no such member, exactly because it cannot promise termination.}
\end{solution}

\begin{solution}{iq:iv3:cpp:22}
It prints \texttt{2 2}, and \texttt{freed} never prints. Each node is held by a local and by the other node. When the locals go, each count falls to 1, so neither reaches 0. LeakSanitizer reports 64 bytes in 2 allocations. \Cref{lst:iv3:cp:cyclefixed} breaks the cycle with one weak link, and the test shows both destructors running.
\omcode{code/interviews-3/03-cpp/cpp/cycle_fixed.cpp}{1}{22}{One strong link and one weak one.}\label{lst:iv3:cp:cyclefixed}

Reference counting cannot see a cycle, as it only knows local counts, and nothing outside now refers to the pair. A garbage collector that traces from roots would find them, but C++ has none.

The choice is about which link means ownership. Here \texttt{next} owns and \texttt{prev} watches. A \texttt{weak\_ptr} does not keep the object alive and cannot be dereferenced directly, so you must call \texttt{lock()}. That gives a \texttt{shared\_ptr}, or null if the object has gone, and that is why the listing tests it. While you hold the result, the object cannot die under you. That is the count of 2 on the line \texttt{alive 2}.

After the block, \texttt{expired()} is true, and \texttt{lock()} gives null.

\iqwrong{``\texttt{use\_count} tells me whether it is safe to use.'' On several threads it can be stale the moment it returns. \texttt{lock()} is the atomic check.}
\iqfollow{What does a raw pointer offer over a \texttt{weak\_ptr} here?}{It costs nothing, and is fine if the owner surely outlives it. A \texttt{weak\_ptr} is for when it may not.}
\end{solution}

\begin{solution}{iq:iv3:cpp:23}
\Cref{lst:iv3:cp:unique} is one such class, as large as a raw pointer, 8 bytes on x86-64. Self-move is safe, as \texttt{release} empties the object first and \texttt{reset} takes the same pointer back. It behaves like \texttt{std::unique\_ptr} on 20\,000 random operations, with no leak under AddressSanitizer.
\omcode{code/interviews-3/03-cpp/cpp/unique.hpp}{4}{28}{A small owning pointer.}\label{lst:iv3:cp:unique}

Ownership means one owner, so copying is deleted, and a move hands over the pointer and leaves null. Both are \texttt{noexcept}, so a vector of these moves during growth.

\texttt{reset} stores the new pointer first because the destructor of the old object may reach back here. If a node owns a child through a \texttt{Unique}, and that child's destructor ends up calling \texttt{reset} on the same pointer, an old-first version would delete the object twice. \texttt{std::exchange} leaves the member consistent before any outside code runs.

\texttt{explicit} on the constructor matters. Without it, a raw pointer would convert silently, and passing one to a function taking a \texttt{Unique} would delete it at the end of the call.

\texttt{delete nullptr} is defined and does nothing, so no check is needed in the destructor. The type must be complete at that point, or the compiler cannot call its destructor, which matters when the class is used with a forward declaration, as in the pimpl idiom.

\iqwrong{\texttt{delete p\_; p\_ = o.p\_; o.p\_ = nullptr;} in move assignment, with no check. On self-move it deletes the object and ends up null: the value is silently destroyed.}
\iqfollow{What does the real one add?}{A deleter type, array support, and conversion from a pointer to a derived class.}
\end{solution}

\begin{solution}{iq:iv3:cpp:24}
Both print \texttt{destroyed} when the last \texttt{shared\_ptr} goes. With \texttt{new}, there is 1 free then and 2 after the \texttt{weak\_ptr} goes. With \texttt{make\_shared}, there are 0 frees then, and 1 after. So in the second case, the megabyte stays allocated until the last \texttt{weak\_ptr} has gone.

\texttt{make\_shared} puts object and control block in one allocation. The control block must outlive the object, since a \texttt{weak\_ptr} needs its counts to answer \texttt{expired()}. The destructor runs when the strong count reaches 0, but the memory is one block, so it cannot be returned in part. With \texttt{new}, there are two blocks, and the object's goes first.

So avoid \texttt{make\_shared} for a large object watched by long-lived weak pointers, such as a cache holding weak references. A destructor that frees a member's memory does help, as only the object's own bytes stay.

Otherwise, \texttt{make\_shared} is the usual choice. One allocation instead of two means less time and less overhead, and object and counts are close together, which helps the cache. You also do not write \texttt{new}. It cannot take a custom deleter, and it needs a public constructor.

\iqwrong{``The destructor has run, so the memory is free.'' Destruction and deallocation are two steps, and here they happen at different times.}
\end{solution}

\begin{solution}{iq:iv3:cpp:25}
\texttt{a} and \texttt{b} each create their own control block for one object, and each count is 1, so they do not know about each other. Whichever dies second deletes the object again, and AddressSanitizer reports \texttt{attempting double-free}.

Ownership is shared only by copying a \texttt{shared\_ptr}. Constructing from a raw pointer always starts a new count. So keep the raw pointer out of it: \texttt{auto a = std::make\_shared<Order>(); auto b = a;} gives counts of 2.

\texttt{std::shared\_ptr<Order>(this)} starts a second count on an object already owned, which is the same bug, with the two owners even further apart.

The right way, in \cref{lst:iv3:cp:sharethis}, is to derive from \texttt{std::enable\_shared\_from\_this} and call \texttt{shared\_from\_this()}. When a \texttt{shared\_ptr} takes ownership, it stores a \texttt{weak\_ptr} to itself inside that base, and the call locks it, which shares the same control block.
\enlargethispage{-3\baselineskip}%
\omcode{code/interviews-3/03-cpp/cpp/share_this.cpp}{1}{18}{Handing out a pointer to yourself.}\label{lst:iv3:cp:sharethis}

So the object must already be owned by a \texttt{shared\_ptr}. If it is not, as with \texttt{plain}, the weak pointer is empty, and since C++17 the call throws \texttt{std::bad\_weak\_ptr}. It is also unavailable in the constructor, since no \texttt{shared\_ptr} has yet taken ownership.

\iqwrong{``The counts are atomic, so two owners are safe.'' Atomicity is about one count seen by many threads. There are two counts here.}
\end{solution}

\begin{solution}{iq:iv3:cpp:26}
\texttt{push\_back} needs more capacity than the vector's 3, so it moves the elements to a new block and frees the old one. \texttt{first} still points into the freed block. AddressSanitizer reports \texttt{heap-use-after-free} at the line \texttt{first = 9}.

Any operation that reallocates invalidates all references, pointers and iterators into a vector. If the capacity is not exceeded, those before the insertion point stay valid. So this bug depends on the capacity at the time, which is why it hides: a vector that already had spare room runs correctly.

Two fixes. Call \texttt{v.reserve(4)} first, and then the reference stays valid, as the test confirms with a clean run under the sanitizer. That is fragile, since the reserve has to cover every later insertion. Or hold an index, \texttt{std::size\_t first = 0}, and write \texttt{v[first] = 9}, which survives any reallocation, and is the usual answer.

\texttt{std::deque} has different rules. \texttt{push\_back} and \texttt{push\_front} invalidate iterators but not references to elements, so this exact code is correct on a deque. An insertion in the middle invalidates both.

\iqwrong{``It printed 1 here, so the write was lost, and that is the only harm.'' Writing to freed memory can corrupt another object or the allocator's own data, far from this line.}
\iqfollow{What about a range-for that pushes onto the vector it walks?}{Same bug: its iterators are invalidated by the first reallocation.}
\end{solution}

\begin{solution}{iq:iv3:cpp:27}
It prints \texttt{0 1 1}. \texttt{auto y = w[0]} copies an \texttt{int}. \texttt{vector<bool>} stores bits, and a bit has no address, so \texttt{v[0]} returns a small proxy object, and \texttt{auto x} is that proxy. Assigning to \texttt{x} writes through to the bit, so \texttt{v[0]} becomes true.

\texttt{bool\& b = v[0];} does not compile, since the proxy is not a \texttt{bool}. Neither does \texttt{auto\& b = v[0];}, because the proxy is a temporary. \texttt{bool b = v[0];} compiles, and it makes an independent copy, which is what \texttt{auto} did not do.

Two different indices of \texttt{vector<bool>} can share a byte, so writing them means reading and writing that byte, and the standard does not make this safe: \texttt{vector<bool>} is the named exception to the rule that different elements of a container can be written concurrently. ThreadSanitizer reports a data race in the test for \texttt{vector<bool>}, and none for \texttt{vector<char>}.

So \texttt{vector<bool>} is not really a container of \texttt{bool}: \texttt{bool*} cannot point at an element, and generic code that takes a \texttt{T\&} fails. Use \texttt{vector<char>} when you want the usual behaviour, and \texttt{std::bitset} or \texttt{vector<bool>} when you want the memory saving of one bit each.

\iqwrong{``It saves memory, so it is always better.'' Each access takes a shift and a mask, and the race rule is different. On a tight loop, bytes can be faster.}
\end{solution}

\begin{solution}{iq:iv3:cpp:28}
It prints \texttt{1 1}, \texttt{1} and \texttt{0}. The standard promises the second number of the first line and the second line: references and pointers into \texttt{unordered\_map} and \texttt{map} stay valid across insertions. The first number says that the table rehashed, which depends on the starting bucket count that the library chooses. The \texttt{0} says that the vector moved, and no promise covers that value either.

A rehash invalidates iterators of an \texttt{unordered\_map}, and not references or pointers to its elements. The elements live in separately allocated nodes in libstdc++, and the standard's rule forces such a design. That is why \texttt{unordered\_map} cannot be a flat open-addressed table, which is the price of this promise.

For pointers to orders held in a container: \texttt{std::map}, \texttt{std::unordered\_map}, \texttt{std::list} and \texttt{std::deque} (at the ends) all allow it. \texttt{std::vector} does not, unless the capacity is fixed in advance. So a book often keeps orders in a pool with stable addresses, with an index on top, or stores integer indices, which also make better use of the cache.

Erasing is a separate matter. It invalidates references only to the erased element, on \texttt{map} and \texttt{unordered\_map}.

\iqwrong{``An \texttt{unordered\_map} moves everything on rehash, so pointers die.'' That confuses it with a vector. Its nodes stay put, and only the bucket array is rebuilt.}
\iqfollow{What does a flat hash table give up to run faster?}{Exactly this: its elements move on growth, so pointers do not survive.}
\end{solution}

\begin{solution}{iq:iv3:cpp:29}
It prints \texttt{1}, \texttt{1} and \texttt{1000000}. The first count comes from \texttt{push\_back(T(1))}, which builds a temporary and moves it in. \texttt{emplace\_back(2)} builds in place, with no move, so the count stays at 1. The \texttt{push\_back} of a plain 1000000 does not compile, and that is the point.

\texttt{emplace\_back} helps when you pass constructor arguments and a move is costly, or not available at all. If you already have an object, \texttt{push\_back(std::move(x))} and \texttt{emplace\_back(std::move(x))} do the same work. For a cheap type such as a small struct of integers, the difference is near nothing once optimised.

The danger is that \texttt{emplace\_back} can call \texttt{explicit} constructors. The vector constructor taking a size is \texttt{explicit} so that a number cannot quietly become a vector. So the \texttt{emplace\_back} in the listing makes a row of a million zeros, where the author may have meant a row holding one number. \texttt{push\_back} needs an implicit conversion, and there is none, so the mistake is caught. With braces, \texttt{push\_back(\{1000000\})} gives the row of one element.

So use \texttt{push\_back} by default, for its checking, and \texttt{emplace\_back} when you are really building in place.

\iqwrong{``\texttt{emplace\_back} is always faster, so use it everywhere.'' It can be equal, it can be slower to compile, and it switches off a safety check.}
\end{solution}

\begin{solution}{iq:iv3:cpp:30}
AddressSanitizer reports a \texttt{heap-buffer-overflow}, a read past the end, in \texttt{\_\_unguarded\_partition}, inside the sort. The comparator broke the requirement of a strict weak ordering: \texttt{comp(a, a)} must be false, and with \texttt{<=} it is true. Breaking it makes the call undefined, so this crash is one allowed result among many.

The implementation depends on the promise. libstdc++'s partition step scans forward \emph{without a bounds check}, on the guarantee that median-of-three left an element in the range that is not less than the pivot. With \texttt{<=}, an equal element also counts as less, so the scan does not stop there and walks off the end of the array.

It passes on small inputs, as libstdc++ uses insertion sort for short ranges, with no such step: the threshold is 16 elements in this version. It passes on distinct values too, because \texttt{<=} and \texttt{<} then give the same answers. So only many equal keys show it, which is what happens the first time a real data set arrives with repeated prices.

For two fields, compare a tuple: \texttt{std::tie(a.px, a.t) < std::tie(b.px, b.t)}, which gets the logic right, or in C++20 a defaulted \texttt{operator<=>}. Writing it by hand invites the mistake: \texttt{a.px < b.px || a.t < b.t} is not a strict weak ordering either.

\iqwrong{``It sorts a little oddly, but the worst case is a wrong order.'' A broken comparator can read and write outside the array.}
\iqfollow{What other things need this?}{\texttt{std::map}, \texttt{std::set} and heaps rely on the same rule.}
\end{solution}

\begin{solution}{iq:iv3:cpp:31}
It prints \texttt{0 1}: the first check misses the overflow, and the built-in catches it. Signed overflow is undefined, so the compiler may assume that \texttt{a + b} does not overflow. For that case, \texttt{a + b < a} is the same as \texttt{b < 0}, which is what it compiles.

With $b = 1$, that test is false, whatever $a$ is, so the check can never fire. The code reads as if overflow wraps, and on the hardware it does, but the optimiser works from the language's rules, not from the machine's.

Under UBSan, the run reports \texttt{signed integer overflow: 2147483647 + 1} and prints \texttt{1 1}, since the sanitizer's build does the addition and it wraps. With \texttt{-fwrapv}, overflow is defined to wrap, so the check works, and the program prints \texttt{1 1} here too. Both change what is being asked, and neither makes the code correct C++.

The portable check tests before adding: for $b > 0$, \texttt{a > INT\_MAX - b}, and for $b < 0$, \texttt{a < INT\_MIN - b}. Neither side can overflow. \texttt{\_\_builtin\_add\_overflow} is a compiler extension in g++ and Clang that does the same, and here it reads the CPU's overflow flag with a \texttt{seto}, as our test confirms.

\iqwrong{``It must be a compiler bug, since the hardware wraps.'' The program has the bug. g++ 11.4 applies this rewrite even at \texttt{-O0}: the listing prints \texttt{0 1} there too.}
\end{solution}

\begin{solution}{iq:iv3:cpp:32}
It prints \texttt{-2147483648}, \texttt{2}, \texttt{-2147483648} and \texttt{4294967295}. The first and last are defined by the standard, and the second and third are undefined, so UBSan flags those two.

\texttt{1 <{}< 31} is defined since C++20, as signed integers are two's complement and the bits go where they go, which makes this \texttt{INT\_MIN}. Before C++20, its value was implementation-defined. \texttt{u - 1} on an unsigned 0 wraps to $2^{32} - 1$, which is always defined.

\texttt{1 <{}< 33} on a 32-bit type has a shift count at or above the width, which remains undefined in C++20. It prints 2 because x86's shift instruction uses only the low 5 bits of the count, so 33 acts as 1. That is this machine, and nothing obliges another machine or another compiler to agree.

\texttt{std::abs(INT\_MIN)} has no answer in \texttt{int}, as $2^{31}$ does not fit, so it is undefined. On this machine the negation wraps back, which leaves a negative absolute value. Any code that assumes an absolute value is non-negative then fails, for instance by using it as an index modulo a table size.

UBSan's reports read \texttt{shift exponent 33 is too large} and \texttt{negation of -2147483648}.

\iqwrong{``Undefined means crash.'' Here both undefined lines print plausible numbers. That is the usual way, which is why you need the sanitizer.}
\iqfollow{And \texttt{INT\_MIN / -1}?}{Also undefined, as the true answer does not fit, and on x86-64 the division instruction traps. Our test sees \texttt{SIGFPE}.}
\end{solution}

\begin{solution}{iq:iv3:cpp:33}
\texttt{one} and \texttt{two} are safe, and \texttt{three} dangles. Binding a temporary directly to a \texttt{const} reference extends the temporary's life to that of the reference, and this also covers a member of the temporary, which is \texttt{two}. AddressSanitizer reports \texttt{stack-use-after-scope} at the last read.

The extension is a property of the initialisation, and the compiler applies it where it sees the temporary being bound. It does not look inside a function. For \texttt{same(make())}, the temporary binds to the parameter \texttt{s}. That lasts until the end of the full expression, and the function returns a reference to it. The compiler cannot know, in general, that the returned reference refers to that argument, so \texttt{three} is bound to an object that is about to die.

The same goes for anything that returns a reference to its argument, which includes \texttt{std::max}, \texttt{std::min} and \texttt{std::move}. A member function that returns a reference to a member of a temporary is the same: \texttt{Holder\{make()\}.get\_name()} would dangle, where the direct \texttt{.name} did not.

So when in doubt, take a value. \texttt{std::string three = same(make());} copies while the temporary is still alive, and costs one copy.

\iqwrong{``A \texttt{const} reference always keeps a temporary alive.'' Only one bound to it directly, in the declaration.}
\iqfollow{Is the range-for form safe?}{\texttt{for (auto\& x : make\_vector())} is, as the range is bound directly. A chain such as \texttt{make\_holder().items()} is not, before C++23.}
\end{solution}

\begin{solution}{iq:iv3:cpp:34}
It prints \texttt{Derived 1}, \texttt{Derived 2}, \texttt{Base::g(int)} and \texttt{Derived::g(double)}. Default arguments are chosen at compile time from the static type, and the function body at run time from the object. Names in a derived class hide all base-class overloads of that name.

So \texttt{b.f()} takes its default from \texttt{Base}, as \texttt{b} is declared a \texttt{Base\&}, and it then dispatches to \texttt{Derived::f}, which gets 1. Two defaults for one function are a trap with no warning. The fix is not to redefine defaults on an override. Better, put the default in one non-virtual function that calls a virtual function with no defaults.

Lookup for \texttt{d.g(1)} stops at the first scope that holds the name, which is \texttt{Derived}. \texttt{Base::g(int)} is not even considered, so 1 converts to \texttt{double}. The fix is \texttt{using Base::g;} in \texttt{Derived}, which brings the base overloads into the set, and then \texttt{d.g(1)} picks the exact match.

A compiler option can warn about one case: \texttt{-Woverloaded-virtual} reports a function hiding a virtual function of the base. Here \texttt{g} is not virtual, so that option does not help, and the listing compiles clean. So review for it.

\iqwrong{``\texttt{override} would have caught this.'' It checks that \texttt{f} overrides something. It says nothing of defaults, and \texttt{g} is not virtual.}
\end{solution}

\begin{solution}{iq:iv3:cpp:35}
It prints \texttt{1 8 4 8 4}, \texttt{8 8 16} and \texttt{16 32 32 24}. An object must have a unique address, so two elements of an array of \texttt{Empty} need different addresses, and the size must be at least 1. \texttt{Base1} is 4 and \texttt{Member} is 8 by the empty base optimisation.

An empty base can share an address with the first member, as a base subobject need not be distinct in that way. An empty member, however, takes a byte, and padding rounds \texttt{Member} up to 8. C++20's \texttt{[[no\_unique\_address]]} gives the member the same treatment, which is why \texttt{Attr} is 4. \texttt{Virt} is 8 for its pointer to a virtual table.

A lambda with no captures has an empty type, and its call is known from the type. libstdc++'s \texttt{unique\_ptr} stores the deleter so that it occupies nothing, so the size stays 8. A function pointer deleter is a real value, chosen at run time, and must be stored, which gives 16.

\texttt{shared\_ptr} holds two pointers: one to the object, and one to the control block with the counts. Two are needed because the object pointer may differ from the owned object, for a base class or with aliasing.

\texttt{std::function} is 32 here, since it holds a small buffer for the callable. \texttt{std::string} is 32, with room for 15 characters, and a vector holds three pointers. The standard fixes none of these sizes.

\iqwrong{``An empty struct takes no space.'' As a base, usually so. As a member or an array element it does, unless told otherwise.}
\end{solution}

\begin{solution}{iq:iv3:cpp:36}
It fails at link time, with \texttt{undefined reference} to \texttt{total<int>}. Each file is compiled alone. \texttt{main.cpp} sees only the declaration, so it assumes that someone else provides \texttt{total<int>}. \texttt{sum.cpp} sees the body, but nothing there asks for any \texttt{T}, so it generates no code.

A template is a recipe, and it produces code only when instantiated with a type, and only where the compiler can see both the request and the recipe. No single file here has both.

Two fixes. The usual one is to put the body in the header, so that every caller can instantiate it. A template definition may appear in several files, and the linker keeps one copy.

The second is explicit instantiation. Leave the body in \texttt{sum.cpp} and add the line \texttt{template int total(const std::vector<int>\&);} there, which makes it generate that one version. This keeps headers light and speeds up builds, and our test builds and links with it.

The cost is the closed list. A caller with \texttt{std::vector<long>} gets the same link error, as the test shows, until someone adds a second line to \texttt{sum.cpp}. A user of a compiled library cannot add one at all.

\iqwrong{``\texttt{export} solves this.'' That keyword's template meaning was removed in C++11. C++20 modules are the modern route, and they need build-system support.}
\iqfollow{What does the header form cost?}{Build time: each file that includes it parses the body and may instantiate it, and \texttt{extern template} can reduce that.}
\end{solution}

\begin{solution}{iq:iv3:cpp:37}
It prints \texttt{integer real integer}, then \texttt{10}. Each \texttt{kind} is viable only if its concept holds for the argument type, and so exactly one survives. \texttt{notional(1, 2)} fails with an error at the call, saying that the constraint \texttt{Priced<int>} is not satisfied, and which requirement failed.

A concept is a named compile-time predicate over types. \texttt{requires(const T\& t) \{ \dots \}} tests whether the expressions in it would compile, and it can check the result type too. So \texttt{Priced} asks that \texttt{t.price()} be valid and convert to \texttt{double}. Here, \texttt{int} has no member \texttt{price}. If two constrained overloads both apply, the more constrained one is chosen.

Without a constraint, the template would be accepted for \texttt{int}, and the error would come from inside its body, on the line \texttt{t.price() * qty}. In a deep library, that message can run to hundreds of lines, and points nowhere near the mistake.

Before C++20, the usual means was SFINAE: an extra template parameter defaulted to \texttt{std::enable\_if\_t<std::is\_integral\_v<T>{}>} removes the overload silently when deduction fails. It works and reads poorly, and two overloads differing only by such a default argument are a redefinition, so you need a dummy parameter to separate them. \texttt{if constexpr} covers some cases more simply, from C++17.

\iqwrong{``A concept is a kind of base class.'' It has no objects and adds no virtual calls, and a type needs no declaration to satisfy one.}
\end{solution}

\begin{solution}{iq:iv3:cpp:38}
It prints \texttt{300 2 24}, \texttt{1}, \texttt{bad access} and \texttt{0 1}. A variant holds its alternative in place, with no allocation: 24 bytes here, for the largest alternative of 16 bytes aligned to 8, plus the index. \texttt{get\_if} returns null for the wrong type, and \texttt{get} throws.

Compared with a base class and three derived classes, there is no heap and no pointer to chase, so a vector of messages is contiguous. \texttt{std::visit} makes the compiler check that every type is handled, since \texttt{Overload} fails to compile if one is missing. So adding a type breaks every visitor until it is handled, which is either what you want or a cost, depending on who owns the code. Adding an operation is easy, as it is a new visitor. With inheritance it is the other way round: a new type is easy and a new operation touches every class. A closed set of known messages suits the variant, and an open set, say plug-in strategies, suits inheritance.

\texttt{std::visit} dispatches on the index, which need not be cheaper than a virtual call, so the gain to claim is in memory and layout. We timed nothing here.

Yes, it can hold nothing. Changing the type destroys the old value before building the new one, and if that throws, nothing remains. That is \texttt{valueless\_by\_exception()}, and \texttt{index()} returns \texttt{variant\_npos}, as the test shows when \texttt{Fragile} has a \texttt{std::string} member. libstdc++ avoids this for a small trivially copyable type, by building a temporary first: with an \texttt{int} member in place of the string, the same program prints \texttt{0 0}.

\iqwrong{``A variant cannot be empty.'' It can, only in that one failure state, and \texttt{std::visit} on it throws.}
\end{solution}

\begin{solution}{iq:iv3:cpp:39}
You must supply a hash and an equality test, and without a hash the map does not compile. The largest bucket holds 50 entries under the weak hash and, in the run here, 4 under the mixed one. The weak hash gives correct answers, and costs time: every lookup walks a chain of up to 50 keys.

Equality is needed because different keys can share a hash, and the map must tell them apart. C++20 can generate it with \texttt{= default}. A hash may be passed as a template argument, as here, or as a specialisation of \texttt{std::hash<Key>}.

\texttt{Weak} uses the symbol only, so all 50 venues of a symbol have equal hashes, and fall in one bucket whatever the table's size. The map still finds the right entry by comparing keys. So a broken hash is a performance bug and never a correctness bug, as long as equal keys hash the same. That is the one hard rule: if two equal keys hashed differently, lookups would fail.

\texttt{Mixed} packs both fields into 64 bits, and multiplies by an odd constant. libstdc++ takes the hash modulo a prime bucket count, so it would do tolerably on the packed value alone, but a table reduced with a mask would use only low bits, and multiplying does not help there, since it moves information only upward: the low bits of the product still depend only on the venue. Such a table must take the top bits of the product, or first xor the high bits down.

The worst case in the limit is adversarial: if an attacker can choose keys with one hash, every operation becomes linear.

\iqwrong{\texttt{h(symbol) \^{} h(venue)}. XOR is symmetric, and equal fields cancel, so (1, 2) and (2, 1) collide, and (5, 5) hashes to 0.}
\end{solution}

\begin{solution}{iq:iv3:cpp:40}
It prints \texttt{2 1}, \texttt{2 3}, \texttt{2} and \texttt{False True}. Boxing is copying a value type into a new object on the heap, so it can be held through \texttt{object} or an interface. C\# does it in \texttt{object box = a}, in \texttt{ICounter c = a}, and for each of \texttt{o1} and \texttt{o2}.

So \texttt{box} holds a copy with \texttt{Qty} 1, and the later change to \texttt{a} does not reach it. \texttt{c.Add()} changes the box, so \texttt{a} stays 2, while the box reads 3. The trap is the method looking like it changes \texttt{a}. Reading a list element gives a copy too, so the third line shows the list unchanged. And two boxes of 5 are two objects, which is the last line: not the same reference, but equal values.

Each box is also an allocation for the collector to clean up, which is why a hot path in C\# avoids boxing, by using generics, where \texttt{List<Fill>} stores the values inline, and constrained calls.

C++ has no boxing, as it has no common base type and no implicit heap object. An object is a value unless you ask otherwise, a copy is always visible as a copy, and reference semantics are explicit: \texttt{Fill\& ref} or a pointer. \texttt{list[0]} returns a reference, so there, in contrast to C\#, writing through it changes the element. When C++ needs to hold different types behind one interface, it says so with a pointer, a \texttt{std::any} or a \texttt{std::variant}.

\iqwrong{``A C\# \texttt{struct} and a C++ \texttt{struct} are the same idea.'' A C++ \texttt{struct} differs from a \texttt{class} only by default access. In C\#, the keyword picks value or reference semantics.}
\end{solution}


\begin{solution}{iq:iv3:cpp:41}
\Cref{lst:iv3:cp:growth} prints \texttt{1 2 4 8 16 32 64 128 256 512 1024} and then \texttt{1023}. For 1000 insertions, doubling makes 1023 moves, and adding 16 each time makes 31\,248, and that gap widens with $n$. \texttt{v.push\_back(v[0])} in the naive vector reads freed memory, which AddressSanitizer reports as \texttt{heap-use-after-free}.

Doubling moves $1 + 2 + \dots + 512 = 1023$ elements, about one per insertion, so each insertion costs a constant number of moves on average. A fixed increment $k$ moves $k + 2k + \dots$, about $n^2 / 2k$ in all, which is quadratic: our count roughly quadruples from 1000 to 2000 insertions. The increment saves memory and loses time.

A factor $f$ makes about $n / (f - 1)$ moves, so 1.5 makes twice as many as 2, and a factor near 1 becomes the fixed-increment case. The gain of a smaller factor is reuse. With doubling, each new block is larger than all the freed ones together. With 1.5, the blocks are $1, 1.5, 2.25, 3.38, 5.06, 7.59$, and the sixth fits in the space of the first four, which total 8.13, if the allocator joins neighbouring free blocks. Below the golden ratio, this reuse eventually becomes possible.

In the naive code, \texttt{x} refers into the old block, and \texttt{grow} deletes that block before \texttt{x} is read. So the real one builds the new element in the new block first, and only then moves the old ones and frees them.

\iqwrong{``Growth is $O(n)$ per insertion, as elements are copied.'' That is the worst single step. The average over all insertions is constant.}
\iqfollow{What else is wrong with \texttt{new T[c]} here?}{It needs a default constructor, and builds $c$ objects that nobody asked for.}
\end{solution}

\begin{solution}{iq:iv3:cpp:42}
A member \texttt{T value} would have to be built whenever the optional is, which needs a default constructor and an object that nobody asked for. So the storage is raw bytes of \texttt{T}'s size and alignment, and the object appears through placement \texttt{new}. \Cref{lst:iv3:cp:opt} sets the flag only after the constructor succeeds. The misaligned program is not correct, even though it runs.
\omcode{code/interviews-3/03-cpp/cpp/opt.hpp}{5}{39}{A small optional.}\label{lst:iv3:cp:opt}

\texttt{alignas(T)} matters as an array of \texttt{unsigned char} has alignment 1 by itself. If the flag were set first and the constructor threw, the destructor would destroy an object that never existed. A test with that order fails, and the listing's order passes against \texttt{std::optional} on 20\,000 random operations, under AddressSanitizer, with a throwing constructor among them and a 64-byte-aligned type.

\texttt{std::launder} is there since a pointer cast from the byte array does not by itself point at the object made inside it. Assignment takes its argument by value, so self-assignment is safe, as the copy exists before \texttt{reset}.

If \texttt{emplace} throws, the old value is already gone and the optional is empty, which is valid state. That is the basic guarantee, and not the strong one.

As for the misaligned \texttt{double}: UBSan reports \texttt{store to misaligned address} and \texttt{load of misaligned address}, and the program still prints 1.5. An object must be at a suitably aligned address, so this is undefined. x86-64 tolerates plain moves, but the compiler may vectorise the access with instructions that assume alignment, and other architectures may trap.

\iqwrong{``\texttt{char buf[sizeof(T)]} is enough.'' It is the right size at the wrong alignment.}
\end{solution}

\begin{solution}{iq:iv3:cpp:43}
Keep the free list inside the free slots themselves: each one is a union of a next pointer and the bytes of a \texttt{T}, so a free slot's memory holds the link, and no extra memory is needed. \Cref{lst:iv3:cp:pool} does that. \texttt{create} pops the head, and \texttt{destroy} pushes it back.
\omcode{code/interviews-3/03-cpp/cpp/pool.hpp}{6}{42}{A fixed pool with an intrusive free list.}\label{lst:iv3:cp:pool}

There is a trap in \texttt{create}. The new object is built over the bytes that hold \texttt{next}, so \texttt{next} must be read before the constructor runs. Our test builds the version that reads it afterwards: the \texttt{id} field overwrites the link with 0, and the pool reports itself full with 15 slots still free.

If the constructor throws, nothing was created, and the slot goes back on the list, so an exception loses no capacity. The test throws from one creation in about 8, and the live count and the exhaustion test still agree with a model after 100\,000 random operations.

It makes no call to \texttt{operator new} in that run, and each slot is aligned for \texttt{T}, which the test checks with a 64-byte type. A slot is at least the size of a pointer, however small \texttt{T} is. The most recently freed slot comes back first, and is the one most likely to still be in the cache.

It does not catch these: destroying a pointer twice, which corrupts the list with a cycle; using an object after \texttt{destroy}; passing a pointer from another pool. And it is for one thread. A second thread needs a lock or a lock-free list, where popping meets the ABA problem. A debug build could keep a bitmap of live slots to catch the first of these.

\iqwrong{``\texttt{T* p = new T[N]} and a stack of pointers.'' That constructs all $N$ objects at the start and needs a default constructor, so \texttt{create} would have nothing left to construct.}
\end{solution}

\begin{solution}{iq:iv3:cpp:44}
An allocator needs only \texttt{value\_type}, \texttt{allocate}, \texttt{deallocate}, equality, and a converting constructor from the same allocator for another type. \Cref{lst:iv3:cp:arena} is the arena and that allocator. \texttt{deallocate} does nothing, as a bump arena frees everything at once, by dropping the buffer.
\omcode{code/interviews-3/03-cpp/cpp/arena_alloc.hpp}{6}{35}{A bump arena and an allocator over it.}\label{lst:iv3:cp:arena}

\texttt{std::align} advances the pointer past any padding, and fails when the rest of the buffer is too small, at which point we throw. The test takes 1 byte and then an aligned \texttt{double}, and gets an aligned address.

The converting constructor exists because of what a container allocates. You name \texttt{pair<const int, double>}, but a \texttt{std::map} allocates tree nodes, a private type that you never see. It rebinds your allocator to that node type and constructs one from yours, which the constructor makes possible. Equality says whether storage from one can be released through the other, so here it compares the arena.

The allocator holds a pointer to the arena and not the arena itself, as allocators get copied, and all copies must draw on one buffer.

The old blocks are simply lost. A vector growing to 100 ints takes 1, 2, 4 up to 128, which is 255 ints of arena for 100 elements. So call \texttt{reserve} when you know the size, or use a pool that reuses memory.

The test fills a vector and a map from the buffer, makes no call to global \texttt{operator new}, gets the right sums back, and sees \texttt{bad\_alloc} when a request is too large. The buffer must outlive all containers using it, which the type does not enforce.

\iqwrong{``Containers call \texttt{allocate} for the elements I named, so one type is enough.'' That is true of a vector, not of a node-based container.}
\end{solution}

\begin{solution}{iq:iv3:cpp:45}
It prints \texttt{1020 400}, then \texttt{1 0}, then \texttt{bad\_alloc}. Without \texttt{reserve}, the vector asks for 1, 2, 4 up to 128 ints, which is 255 ints or 1020 bytes, and the arena never reuses any. With \texttt{reserve(100)}, it asks once, for 400 bytes.

The second line shows that a string's resource is the buffer, and that global \texttt{new} was called zero times. A \texttt{pmr::vector} constructing a \texttt{pmr::string} passes its own allocator down, so each element uses the same resource. This is called uses-allocator construction, and it is the reason nested containers are so easy here.

The last line comes from the resource's upstream. When the 64-byte buffer runs out, the monotonic resource asks upstream for more, and \texttt{null\_memory\_resource} always throws. That turns ``this path must not reach the heap'' into a checkable failure. With the default upstream, it would quietly go to the heap instead.

The gain over a template allocator is in the type. \texttt{std::vector<int, ArenaAlloc<int>{}>} and \texttt{std::vector<int>} are different types, so a function must be a template to take both. Every \texttt{pmr::vector<int>} is one type, whatever its resource, so the choice is made at run time.

The cost is a pointer in each container, and a virtual call at each allocation, which matters little when allocations are rare. There is also a trap. Copy-constructing a \texttt{pmr::vector} does not inherit the source's resource: the copy uses the default one, which our test shows, so a copy can reach the heap unnoticed.

\iqwrong{``The vector frees its old blocks into the buffer.'' \texttt{do\_deallocate} on a monotonic resource does nothing, and memory only returns when the resource is released or destroyed.}
\end{solution}

\begin{solution}{iq:iv3:cpp:46}
\texttt{a = a} silently wipes the data, and the program prints 0 rather than 7. \texttt{a = b} with a throwing \texttt{get} leaves \texttt{d\_} pointing at freed memory, so the destructor frees it again, and AddressSanitizer reports \texttt{attempting double-free}. \Cref{lst:iv3:cp:assigngood} is correct in both cases, through copy and swap.
\enlargethispage{-3\baselineskip}%
\omcode{code/interviews-3/03-cpp/cpp/assign_good.cpp}{12}{30}{Copy, then swap.}\label{lst:iv3:cp:assigngood}

On self-assignment, \texttt{o} and \texttt{*this} are one object. \texttt{delete[] d\_} frees the data, then \texttt{d\_ = get(n\_)} installs a fresh zeroed block, and \texttt{std::copy} copies that block onto itself. No freed memory is read, so AddressSanitizer stays silent, and the data is just gone. That self-copy also breaks \texttt{std::copy}'s precondition, so formally it is undefined; on this library it is a harmless \texttt{memmove}. An \texttt{if (this != \&o)} would fix this fault and leave the other in place.

The second fault is the order: the old state is destroyed before the new one is known to exist. If \texttt{get} throws, \texttt{d\_} still holds the deleted pointer.

In the new version, the copy is made first, while \texttt{*this} is untouched, and that is the only line that can throw. The swap cannot, and the destructor of \texttt{copy} frees the old data. Self-assignment needs no check. The test prints \texttt{7}, \texttt{caught 7} and \texttt{8}: the value survives self-assignment, and a failed assignment leaves it unchanged.

The cost is an allocation on every assignment, even when the existing block is big enough. A \texttt{std::vector} avoids that by copying into its existing capacity, and so gives up the strong guarantee for its copy assignment: if an element's assignment throws partway, the vector holds a mix. On a hot path, that is the right trade, provided the caller can live with it.

\iqwrong{``Add \texttt{if (this == \&o) return *this;} and it is fixed.'' That handles a case that barely happens, and leaves the failure that does.}
\end{solution}

\begin{solution}{iq:iv3:cpp:47}
It prints \texttt{1}, then \texttt{3f800000}. An \texttt{int*} and a \texttt{float*} are assumed not to point at the same object, so the compiler returns 1 without reading \texttt{*a} again. With \texttt{-fno-strict-aliasing}, it makes no such assumption, reloads, and the first line becomes 0.

The rule: an object may be read only through its own type, a signed or unsigned version of it, or a \texttt{char}, \texttt{unsigned char} or \texttt{std::byte}. Reading an \texttt{int} through a \texttt{float} is undefined, so \texttt{main} is at fault here, and not the optimiser. Assumptions like this one let the compiler keep values in registers across unrelated stores, so a loop writing through a \texttt{float*} need not reload an \texttt{int} each time round.

The second line shows the legal way to see a float's bits: \texttt{std::bit\_cast} gives a new value with the same bits, here 0x3f800000 for 1.0.

For a packet, \texttt{*reinterpret\_cast<const std::uint32\_t*>(p)} breaks two rules: aliasing, and alignment, since a field in a packet may sit at any offset. The right way is to declare a \texttt{std::uint32\_t v} and \texttt{std::memcpy(\&v, p, 4)}. It looks slow, and costs nothing here: at \texttt{-O2}, our test finds a single \texttt{movl (\%rdi), \%eax}, and no call. Convert the byte order afterwards, if the wire format needs it.

Many large code bases compile with \texttt{-fno-strict-aliasing} as a safety net, which is a choice, and leaves the program non-portable.

\iqwrong{``A \texttt{union} of the two is the legal way.'' That is allowed in C, and it is undefined in C++, though g++ documents it as an extension.}
\iqfollow{And casting a struct over the buffer?}{The same two rules are broken, and the layout has padding to think about too. Copy the fields, or \texttt{memcpy} into a struct whose layout you have checked.}
\end{solution}

\begin{solution}{iq:iv3:cpp:48}
\texttt{read\_bad} compiles to one load and a return, and has no check at all. \texttt{read\_good} tests the pointer, and has a branch to return $-1$. The dereference on the first line tells the compiler that \texttt{p} is not null, because a null dereference is undefined, so \texttt{if (!p)} can never be true, and it is deleted.

The compiler makes no special effort here. It propagates a fact it learned, as it does everywhere, and then removes dead code. Each step is valid, and the result shocks only because the check was the author's safety net. A warning is not promised, and our build with \texttt{-Wall -Wextra} gives none.

On \texttt{read\_bad(nullptr)} on this platform, the load reads address 0, and the process dies with a segmentation fault. So the check was useless even before it was removed, as the crash comes on the line above. However, that is not guaranteed either: if \texttt{v} were unused, the load could vanish too, and the program would carry on in a state its author believed impossible.

\texttt{-fno-delete-null-pointer-checks} keeps the check, and our test finds the \texttt{test} instruction back in the assembly. On a system where address 0 can be mapped, that matters. As a fix for this code, it is the wrong answer: the line order is the bug, and the flag only hides it for one compiler. Move the check to before the use, as \texttt{read\_good} does.

The general lesson is that undefined behaviour reaches backward and forward. An operation in the future can change what the compiler assumes about the past, so ``it only crashes if we get there'' is not reliable.

\iqwrong{``The optimiser removed a check, so that is a compiler bug.'' The standard lets it assume the absence of undefined behaviour.}
\end{solution}

\begin{solution}{iq:iv3:cpp:49}
With no arguments, AddressSanitizer reports \texttt{stack-use-after-scope} at line 6. With three, the program runs clean and prints 4. \texttt{std::max} returns a reference to whichever argument wins. With \texttt{argc} equal to 1, that is the temporary holding 2, which is gone by the end of line 5. With \texttt{argc} equal to 4, it is \texttt{argc} itself, which is still alive.

So the bug depends on the data. A test whose \texttt{argc} is larger than 2 passes, and nothing warns that the other case dangles.

The reference binding \texttt{m} extends no lifetime. The temporary is bound to the function's parameter, and it lives until the end of the full expression. What comes back from the function is a new reference, and the compiler does not know that it may be this temporary.

\texttt{std::max} returns a reference so that it does not force a copy. It works for a type that cannot be copied or is costly to copy, and its result can be used in place, within the same expression. That is safe as long as the result is not kept.

\texttt{std::max(1, 2.0)} does not compile, since \texttt{T} is deduced from both arguments, and comes out as \texttt{int} and as \texttt{double}. No conversion is considered during deduction. You write \texttt{std::max<double>(1, 2.0)}.

The fix here is to hold a value: \texttt{int m = std::max(argc, 2);}. That copies while the temporary is alive, and for an \texttt{int} the copy is nothing. The same holds for \texttt{std::min} and \texttt{std::clamp}.

\iqwrong{``A \texttt{const} reference extends the temporary, so this is safe.'' The extension applies only to a temporary bound directly by that reference.}
\iqfollow{And \texttt{std::max(a, b).size()} on two strings?}{Fine: the call and the use are in one full expression, so nothing has died yet.}
\end{solution}

\begin{solution}{iq:iv3:cpp:50}
Mode \texttt{A} is a data race, and ThreadSanitizer reports it. The count inside the control block is atomic, but the \texttt{shared\_ptr} object is not: it is two pointers, and one thread writes them as another reads them. Mode \texttt{B} goes through \texttt{std::atomic\_store} and \texttt{std::atomic\_load}, and is clean.

What can go wrong is concrete. The reader loads the control-block pointer and is then paused. The writer replaces \texttt{config}, the old count reaches zero, and the block is freed. The reader wakes, increments a count in freed memory, and holds an object that no longer exists. It can also see the new object pointer with the old block, since the two are not read together.

The rule is the usual one for a type: concurrent reads of one object are fine, as is any access to distinct objects, and a write to one object alongside anything else on it is a race. The atomic count is what makes distinct objects safe even when they share an owned object. So four threads with their own copies are not a race, and ThreadSanitizer finds none in \cref{lst:iv3:cp:sharedcopies}.
\enlargethispage{-6\baselineskip}%
\omcode{code/interviews-3/03-cpp/cpp/shared_copies.cpp}{6}{18}{Each thread has its own copy.}\label{lst:iv3:cp:sharedcopies}

\texttt{std::atomic<std::shared\_ptr<int>{}>} is the C++20 way, but libstdc++ 11 does not have it: g++ 11.4 stops at a \texttt{static\_assert}, as \texttt{shared\_ptr} is not trivially copyable. It arrived in GCC 12. So we used the older free functions, which libstdc++ implements by taking one of a small set of mutexes, as its header shows. These functions are deprecated in C++20.

So this path takes a lock, and even a plain copy of a \texttt{shared\_ptr} is an atomic read-modify-write on a line that all readers share. A hot reader would rather keep its own copy, and refresh it now and then.

\iqwrong{``\texttt{shared\_ptr} is thread-safe.'' Its count is. The object holding the two pointers is not.}
\end{solution}

\begin{solution}{iq:iv3:cpp:51}
\Cref{lst:iv3:cp:bq} is one answer. There are two condition variables, so that a push wakes only consumers and a pop only producers, and each wait has a predicate, as a wait can return when the condition does not hold. \texttt{close} sets the flag under the lock and then wakes everyone on both.
\omcode{code/interviews-3/03-cpp/cpp/bq.hpp}{9}{44}{A bounded queue with close.}\label{lst:iv3:cp:bq}

With one condition variable, a push could wake another producer, which finds the queue full and sleeps again, and if \texttt{notify\_one} chose that one, the wake-up is lost, and the waiting consumer sleeps on. Using \texttt{notify\_all} would fix that at the price of many pointless wake-ups.

A wait can return for two reasons without a true condition: a spurious wake-up, which the standard permits, and a stolen one, in which another thread took the item before this thread got the lock. The predicate form loops and checks again.

Setting the flag under the lock is essential. Otherwise a waiter could check its predicate, see no close, and go to sleep just after the notification, and never wake. The notification itself can come after the unlock, so a woken thread does not at once block on the mutex the notifier still holds. Note that \texttt{pop} still returns items after \texttt{close}, so the queue drains, and returns nothing only when it is empty.

The test runs 3 producers and 3 consumers through a queue of capacity 4. It checks the count, 9000, and the exact sum, 27\,009\,000, which would catch a lost or doubled item. That runs under ThreadSanitizer. Then it blocks a producer on a full queue, and confirms that \texttt{close} wakes it with \texttt{false}, and the item already inside can still be popped. A hang would fail the test on its timeout.

\iqwrong{``\texttt{if (q\_.empty()) wait}'' instead of a loop. It returns with an empty queue the first time a wake-up is stolen.}
\end{solution}

\begin{solution}{iq:iv3:cpp:52}
\Cref{lst:iv3:cp:choose} builds the result one factor at a time, cancelling common factors before multiplying. It agrees with an exact computation for all $n \le 75$, reports overflow on exactly the cases that do not fit, and the first overflow appears at $n = 68$, for $k = 31$.
\omcode{code/interviews-3/03-cpp/cpp/choose.cpp}{1}{17}{A binomial coefficient that reports overflow.}\label{lst:iv3:cp:choose}

The factorial formula fails at once, since $21!$ is more than $2^{64}$, though $\binom{21}{10}$ is only 352\,716.

The loop keeps $r = \binom{n - k + i}{i}$ after step $i$. To reach the next value, we need $r \times (n - k + i) / i$, which is an exact division. Multiplying first could overflow when the result fits. So let $g = \gcd(r, i)$. Then $r / g$ and $i / g$ share no factor, so $i / g$ must divide $n - k + i$ exactly. That gives $(r / g) \times ((n - k + i) / (i / g))$, a product of two integers equal to the next binomial, with nothing larger on the way.

Why is no overflow reported falsely? That product is exactly $\binom{n - k + i}{i}$, and these values grow with $i$ up to $k$, since $k \le n / 2$ after taking the minimum. So if a step overflows, then so does $\binom{n}{k}$, which is the last of them.

The smallest $n$ is 68, as $\binom{67}{33}$ is about $1.42 \times 10^{19}$ and fits, and $\binom{68}{31}$ is about $2.19 \times 10^{19}$, above $2^{64} \approx 1.84 \times 10^{19}$. The test checks edge cases too: $k > n$ gives 0, $k = 0$ gives 1, and $n = 2^{64} - 1$ with $k = 1$ returns $n$, and with $k = 2$ overflows.

\texttt{\_\_builtin\_mul\_overflow} is a g++ and Clang extension. A portable form tests \texttt{a > max / b} first.

\iqwrong{\texttt{r = r * (n - k + i) / i} ``as the division is exact''. The division is exact, but the multiplication before it can overflow while the true answer fits.}
\end{solution}

\begin{solution}{iq:iv3:cpp:53}
With \texttt{noexcept} and \texttt{destructor}, the program prints \texttt{terminate called} and the message, and dies with \texttt{SIGABRT}. With \texttt{plain}, it prints \texttt{caught plain}. The \texttt{try} block is never given the first two, as the exception cannot leave the function that promised not to throw.

\texttt{noexcept} is a promise enforced at the function's edge. When the search for a handler reaches the edge of a non-throwing function, the language calls \texttt{std::terminate} rather than carrying on outward. The standard leaves open whether the stack is unwound first, so destructors may not have run.

Destructors are \texttt{noexcept} by default since C++11, without being written so, unless a member's or a base's destructor is declared able to throw. So \texttt{\textasciitilde{}Closer()} is just as non-throwing as \texttt{promised}. The compiler sees this and warns with \texttt{-Wterminate}, which is on by default. The listing turns that off on a line not shown, to make the point.

You can allow it: \texttt{\textasciitilde{}Closer() noexcept(false)}. Our test builds that, and it prints \texttt{caught from destructor}. But if a destructor throws while another exception is already unwinding the stack, there would be two at once, and the program terminates, with no way to catch either. The test shows that too. So a destructor that can run during unwinding must not throw.

A file or socket closer should report failure another way. Offer a \texttt{close()} method that may throw, or return a status for a caller who cares, and let the destructor swallow and log whatever remains.

\iqwrong{``Wrap the code in \texttt{try} and \texttt{catch (.{}.{}.)} in \texttt{main}.'' That is the code shown. Control never returns there.}
\iqfollow{Does \texttt{noexcept} make code faster?}{It removes some work for the caller, and it lets containers move. The second is where it matters most.}
\end{solution}

\begin{solution}{iq:iv3:cpp:54}
It prints 62: \texttt{static\_call} gives $(10 + 1) \times 2 = 22$, and the two virtual calls give 22 and 18. \texttt{Momo} derives from \texttt{Strat<Momo>}, so the base knows its derived type at compile time, and \texttt{static\_cast<Derived*>(this)} is a plain, resolved call. The assembly of \texttt{static\_call} has no \texttt{call} at all, while \texttt{virtual\_call} has \texttt{call *(\%rax)}, an indirect call through the virtual table.

This is the curiously recurring template pattern. A base class template takes the derived class as its parameter. The cast is safe if the object really is a \texttt{Momo}, which holds as long as only \texttt{Momo} derives from \texttt{Strat<Momo>}.

The speed comes mainly from inlining. The compiler can see which function is called, so it copies in the body, and that opens up further optimisation across the call boundary. The indirect call itself is usually cheap when it is predicted well. So the difference depends on what the callee does: for a tiny function in a tight loop it can be large, and for a function that does real work it is small. We measured nothing here, so measure in your own code.

Without \texttt{signal}, the error appears at the line in \texttt{on\_tick}, and only if \texttt{on\_tick} is used, since a member of a class template is instantiated on use. A concept could give a better message.

\texttt{Strat<Momo>} and \texttt{Strat<Other>} are unrelated types, as there is no common base for a container to hold. So CRTP suits a strategy known at compile time. To choose among strategies at run time, use a virtual function, or a \texttt{std::variant} and \texttt{std::visit}. A common mix is one virtual call per message, with statically dispatched code inside it.

\iqwrong{``CRTP is always faster, so remove all virtual functions.'' A compiler can also devirtualise when it knows the type, as with \texttt{final}, and a virtual call that is predicted well costs little.}
\end{solution}

\begin{solution}{iq:iv3:cpp:55}
It prints \texttt{42 2 42 1000}. Only \texttt{int d = twice(argc);} compiles of the four. A \texttt{constexpr} function may run at compile time or at run time, and \texttt{consteval} must run at compile time. A \texttt{constexpr} variable must have a constant initialiser.

So \texttt{twice(21)} initialising a \texttt{constexpr} variable runs in the compiler, and \texttt{twice(argc)} runs in the program, as \texttt{d} needs no constant. \texttt{always(argc)} fails because a \texttt{consteval} call must give a constant, and \texttt{argc} is not one. \texttt{constexpr int f = twice(argc)} fails since the variable demands what the argument cannot give.

The last one is the instructive one. Doubling 2\,000\,000\,000 overflows \texttt{int}. At run time, that would be undefined behaviour, compiled silently. In a constant expression, the standard disqualifies any operation that would be undefined, so the compiler must diagnose it. Constant evaluation is thus a free check for undefined behaviour, for values you can compute at compile time, which is one reason to write a \texttt{static\_assert} on a table.

The \texttt{powers} table is built by a loop inside the compiler, and goes into the program as read-only data. No start-up code is needed, there is no order-of-initialisation issue, and \texttt{powers[3]} just reads a constant. The \texttt{static\_assert} that \texttt{powers[18]} is $10^{18}$ checks the table at every build. Note the guard in the loop, which stops it from computing $10^{19}$: without it, the last multiply overflows, and the build fails.

The cost is compile time, and data in the binary. A table used in a hot path may also compete for cache, which a small computation does not, so a table is not always the faster choice.

\iqwrong{``\texttt{constexpr} on a function means it runs at compile time.'' It only permits that. The call decides.}
\iqfollow{And \texttt{constinit}?}{It requires constant initialisation of a static variable, and the variable can then change, unlike \texttt{constexpr}.}
\end{solution}

\begin{solution}{iq:iv3:cpp:56}
\Cref{lst:iv3:cp:vec1,lst:iv3:cp:vec2} are one implementation. A reallocation goes: allocate the new block, build the new element in it at its final place, transfer the old elements, destroy them and free the old block, and only then update the members. The vector moves when the move cannot throw or when there is no copy, and copies otherwise.
\omcode{code/interviews-3/03-cpp/cpp/vec.hpp}{8}{39}{A vector: storage, copy, move.}\label{lst:iv3:cp:vec1}
\omcode{code/interviews-3/03-cpp/cpp/vec.hpp}{40}{69}{A vector: growth and access.}\label{lst:iv3:cp:vec2}

\texttt{new T[c]} would need a default constructor, and would build $c$ objects, leaving a capacity with live objects in it. Here memory comes from \texttt{std::allocator} as raw storage, with objects created by placement \texttt{new} and destroyed by hand. The size \texttt{n\_} says exactly which slots hold objects, which is the invariant the class rests on.

The order gives both requirements. The new element goes first, while the old block is alive, so an argument that refers into the vector, \texttt{v[0]}, is still valid. Our test does that 40 times with a 40-character string. And until the last steps, nothing visible has changed. If the new element's constructor throws, or a copy during the transfer throws, the catch blocks undo what was done in the new block, and the vector is as it was. \texttt{uninitialized\_copy\_n} destroys its own partial work when it fails.

A throwing move would leave some old elements moved-from, so the transfer copies when it can, as the real vector does. A move-only type has no choice, and the test holds \texttt{unique\_ptr}s.

To test the guarantee, use a type whose copy throws on a chosen count, at each position in turn, and compare the contents and capacity afterwards against a copy taken before. Then check the count of live objects, and run under AddressSanitizer. Here it runs 20\,000 random steps against \texttt{std::vector}. A version made to move always fails it, as the test confirms.

\iqwrong{\texttt{put(); d\_ = fresh;} before the transfer, or the new element built last. Either breaks \texttt{v.emplace\_back(v[0])} or the guarantee.}
\end{solution}

\begin{solution}{iq:iv3:cpp:57}
The increment can be relaxed. The decrement needs \texttt{acq\_rel}. \Cref{lst:iv3:cp:rc} is the class, with count and object in one block, and it matches \texttt{shared\_ptr} on 20\,000 random operations. It is clean under AddressSanitizer and ThreadSanitizer, with four threads racing to drop the last reference.
\omcode{code/interviews-3/03-cpp/cpp/rc.hpp}{5}{33}{A reference-counted pointer.}\label{lst:iv3:cp:rc}

Incrementing is relaxed since a new reference can only be made from an existing one, which keeps the object alive throughout, and the increment need not order anything else.

The decrement has two jobs. A thread that used the object and drops its reference must have all its accesses finished before the count falls, which needs \texttt{release}. And the thread that reaches zero and deletes must see all those accesses as done, which needs \texttt{acquire}. If it is relaxed, another thread's last read can be ordered after the \texttt{delete}, and the destructor may see stale members. A common refinement uses \texttt{release} on the decrement, with an acquire fence only when the count reaches zero, which avoids the acquire on every other decrement.

On x86-64, atomic read-modify-writes are full barriers, so a relaxed version normally runs fine, and a stress test will not fail on the hardware. That is why the test uses ThreadSanitizer, which reasons about the memory model and not the machine. It builds the class with the decrement changed to relaxed and reports a \texttt{data race} on the object's data. The same driver with \texttt{acq\_rel} reports none. So the test can fail the model. We did not run this on a weakly ordered machine.

It does not make one \texttt{Rc} object safe to change from two threads, as that is a data race on \texttt{b\_}, and \texttt{count()} can be stale when it returns. Nor does it protect the \texttt{T}, whose members need their own synchronisation. It has no weak count, no custom deleter and no aliasing, and that is what makes it 8 bytes.

\iqwrong{``Every operation should be \texttt{seq\_cst}, to be safe.'' That is correct, and it costs on weakly ordered machines, and shows you did not work out what each one needs.}
\end{solution}

\begin{solution}{iq:iv3:cpp:58}
\Cref{lst:iv3:cp:fn1,lst:iv3:cp:fn2} are one implementation. You forget the type by writing functions while you still know it: \texttt{Model<F, Small>} generates a table of function pointers, one table per callable type, and \texttt{Fn} keeps a pointer to the table. Beyond the call, you need \texttt{move\_to} and \texttt{destroy}.
\omcode{code/interviews-3/03-cpp/cpp/fn.hpp}{7}{35}{A callable holder: the table for each type.}\label{lst:iv3:cp:fn1}
\enlargethispage{-3\baselineskip}%
\omcode{code/interviews-3/03-cpp/cpp/fn.hpp}{37}{63}{A callable holder: the public side.}\label{lst:iv3:cp:fn2}

Inside the constructor, \texttt{D} is known, and the lambdas in \texttt{Model<D, small>} are written with it. After that, \texttt{Fn} sees only \texttt{void*} and three function pointers. This is type erasure, a virtual table written by hand.

The buffer holds a callable only when it is small enough, not over-aligned, and has a \texttt{noexcept} move. That last condition follows from \texttt{Fn}'s own \texttt{noexcept} move: a callable in the buffer must really move when \texttt{Fn} moves, whereas a heap callable needs only its pointer copied. If a move in the buffer could throw, \texttt{Fn}'s move could not honestly promise otherwise. The promise matters: a growing vector moves a type it cannot copy either way, but if such a move throws halfway, the elements already moved stay moved, and the vector's strong guarantee is gone.

A lambda with four \texttt{long} captures is 32 bytes, and the test counts 0 allocations for it. One with an array of eight is 64 bytes, and costs 1. After being moved from, the source tests false, and the moved callable still works, without copying the 64 bytes. A lambda holding a \texttt{unique\_ptr} works, as we never ask for a copy. Its differential test against \texttt{std::function} runs 20\,000 random steps with a live-object count, which ends at 0, under AddressSanitizer.

\texttt{std::function} has copying, which is the reason it cannot hold a \texttt{unique\_ptr} lambda. It also throws \texttt{bad\_function\_call} when empty, where calling our empty \texttt{Fn} is undefined, and it supports \texttt{target}. C++23 adds \texttt{std::move\_only\_function}, which is closest to this class.

\iqwrong{``Use a virtual base class and \texttt{new} for each callable.'' That works, and it allocates each time. The hand-made table is what allows the small buffer.}
\end{solution}

\begin{solution}{iq:iv3:cpp:59}
\Cref{lst:iv3:cp:var} is one implementation. The storage is the largest size among the types, aligned for the strictest of them, which \texttt{alignas(Ts.{}.{}.)} and \texttt{std::max(\{sizeof(Ts).{}.{}.\})} state directly. The index says what is alive, and the destructor tests each type in a fold. If a constructor throws inside \texttt{emplace}, the variant holds nothing.
\omcode{code/interviews-3/03-cpp/cpp/var.hpp}{8}{47}{A small variant.}\label{lst:iv3:cp:var}

A pack in \texttt{alignas} applies each alignment and keeps the largest, so a 32-byte-aligned \texttt{Wide} forces the whole object to 32, which a \texttt{static\_assert} in the test checks.

The destructor has to run the right destructor for a type known only at run time. The fold expands to one test per type: \texttt{idx\_ == 0 ? destroy first : nothing}, then \texttt{idx\_ == 1} and so on, with \texttt{i++} moving along the pack. \texttt{visit} works just the same. That is a linear chain of tests, which is fine for a few types. \texttt{std::variant} implementations typically generate a table or a switch instead.

\texttt{emplace} destroys the old value first, so that both never share the buffer, and sets \texttt{idx\_} to \texttt{none}. If the constructor then throws, \texttt{idx\_} stays \texttt{none}, and the destructor does nothing. That is the same state as \texttt{valueless\_by\_exception}. Without the \texttt{destroy} call, the old object would never be destroyed, and the test's live count catches that.

\texttt{index\_of} is an immediately invoked \texttt{constexpr} lambda with a fold over the pack: for each type, if it matches, record \texttt{i}, then step \texttt{i}. A comma fold evaluates left to right, so no recursion or helper template is needed. If a type appears twice in the list, the last match wins, where \texttt{std::variant} would refuse the call. A type not in the list must be refused at compile time, which the \texttt{static\_assert} does. Without it, \texttt{emplace<std::string>} on a \texttt{Var<int, char>} compiles and overflows the buffer.

The test runs 30\,000 random steps against \texttt{std::variant} with a \texttt{monostate}, covering a throwing \texttt{emplace}, and compares the index, the visit result and the number of live objects, under AddressSanitizer.

\iqwrong{``Use a \texttt{union} of \texttt{Ts.{}.{}.}.'' A union cannot be declared over a pack directly. It takes a recursive template, which is what the standard library does, and which also allows \texttt{constexpr}.}
\end{solution}

\begin{solution}{iq:iv3:cpp:60}
\texttt{a + b} now returns a small \texttt{Add} object holding references to \texttt{a} and \texttt{b}, and does no arithmetic. The assignment to \texttt{d} runs one loop, asking the whole tree for each element, as in \cref{lst:iv3:cp:expr}. The test counts 0 allocations, where the plain version makes 2. \texttt{a = a + b} is safe here. \texttt{auto e} dangles, and AddressSanitizer reports \texttt{stack-use-after-scope}.
\omcode{code/interviews-3/03-cpp/cpp/expr.hpp}{6}{41}{Expression templates.}\label{lst:iv3:cp:expr}

The type of \texttt{a + b + c} is \texttt{Add<Add<Vector, Vector>, Vector>}, so the whole expression is in the type, and the compiler can inline \texttt{e[i]} down to \texttt{a[i] + b[i] + c[i]}. The \texttt{Expr} concept keeps \texttt{operator+} from applying to arbitrary types.

\texttt{a = a + b} is safe, as each element depends only on the same index, so writing \texttt{a[i]} never affects a later read. An operation that reads other positions is not safe in this way: a matrix product or a shift reads elements that the loop has already written, which is why such libraries make a copy for those cases.

The trap is lifetime. \texttt{a + b + c} creates the inner \texttt{Add} as a temporary, and the outer one stores a reference to it. In \texttt{d = a + b + c}, the temporaries live until the end of the full expression, and that includes the loop. With \texttt{auto e}, the outer \texttt{Add} survives the line, and the inner one does not. Before, \texttt{auto} on a \texttt{Vector} result was harmless, which is why people write this. Nothing fails at compile time.

The repair is to store a named \texttt{Vector} by reference, and an inner expression by value, since an expression is only two references or a nested copy. Our test chooses the member type with \texttt{conditional\_t}, and the \texttt{auto} version then runs clean under AddressSanitizer, giving 6. The named vectors must still outlive \texttt{e}.

Mismatched sizes are not checked here, as \texttt{size} looks only at the left operand. A real library asserts on them. We measured no timing.

\iqwrong{``The compiler would have optimised the temporaries away anyway.'' The allocations in the plain version are real, and the test counts 2.}
\end{solution}
